\section{实验结果与分析}

在本节中，我们展示了基函数优化降低全局刚度矩阵最大特征值的效果，从而节省共轭梯度方法的计算时间。我们的算法使用 C++ 实现，数值测试在配备 24 × AMD Ryzen 9 7900X 12 核处理器和 128GB 内存的笔记本电脑上运行。在所有结果中，共轭梯度的收敛条件为 $\| Ax-b\|_2 / \|b\|_2 < 1e-14$。除基准测试中的模型外，我们使用的模型统计数据列于 \Cref{p2:tab:statistics}。

\begin{table}[t]
  \setlength{\tabcolsep}{4pt}
  \resizebox{0.8\linewidth}{!}{%
    \begin{tabular}{|c|cc|cc|c|}
      \hline
      Example & \#nodes &\#TET & $\gamma_{\max}$ &\#T($\gamma>50$)& opt. time (s)\\
      \hline
      \Cref{fig:cut-simulation} iter-0  & 24 K & 82 K & 1342.61 & 94 & 1.25\\
      \Cref{fig:cut-simulation} iter-5  & 33 K & 103 K & 43781 &702 &1.60\\
      \Cref{fig:cut-simulation} iter-8  & 35 K & 108 K & 43781.4 &979 & 2.14\\
      \Cref{fig:strong-link-demo} & 6K & 20K & 77.8 & 8 & 0.65\\
      \Cref{fig:different-number-of-levels} Nut & 17K & 86K & 124.2 & 28 & 0.31 \\
      \Cref{fig:different-number-of-levels} BCC & 19K & 68K & 220.8 & 89 & 0.48 \\
      \Cref{fig:different-number-of-levels} SphereShell & 87K & 367K & 42 & 0 & 1.53\\
      \hline
    \end{tabular}
  }
  \caption{\textbf{模型统计数据。}我们用 $\gamma$ 表示纵横比。}
  \label{p2:tab:statistics}
  \vspace{-1cm}
\end{table}

\subsection{基准测试}

我们使用 Thingi10k\cite{Thingi10K} 数据集评估了我们的方法，该数据集包含 218 个多样化的 3D 模型。对于每个模型，使用 Tetgen \cite{si_2015_tetgen} 生成四面体网格。我们通过固定模型最长 xyz 边界框维度约 $10\%$ 的范围并施加重力载荷，模拟类似于"悬臂梁"的静力平衡问题。生成的线性系统使用无任何预条件子的共轭梯度方法求解。

为了评估我们方法的有效性，我们收集了应用方法前后最大特征值、收敛所需迭代次数和总时间成本的统计数据。使用 \Cref{eq:elewise_eigmax} 中的准则识别基函数优化的候选单元，参数 q 设为 40。我们还对未使用小波变换处理"强链接"的配置进行了测试。\Cref{fig:benchmark} 中的直方图展示了这些统计指标的降低情况。这些结果清楚地表明，将 \Cref{eq:elewise_eigmax} 中的单元选择准则与"强链接"的小波变换结合使用可产生最显著的性能提升：这种最优方法平均将刚度矩阵的最大特征值降低 $50\%$（最高达 $97\%$），并将共轭梯度求解器的计算时间平均减少 $20\%$（最高达 $87\%$）。

\begin{figure*}[htbp]
  \centering
  \includegraphics[width=\linewidth]{plot_figures/cg_benchmark/integrated-violins.png}
  \caption{Thingi10k 数据集\cite{Thingi10K} 的基准测试，包含 218 个模型。从左到右，分别是最大特征值相对降低 $\frac{\lambda_{\max} - \tilde{\lambda}_{\max}}{\lambda_{\max}}$、共轭梯度方法迭代次数相对降低 $\frac{n - \tilde{n}}{n}$ 和时间成本相对降低 $\frac{t - \tilde{t}}{t}$ 的直方图。从上到下，分别是使用 \Cref{eq:elewise_eigmax} 中基于单元最大特征值排序（无/有小波重排序）和使用 \Cref{eq:rowsum_max} 中基于最大"行和"排序（无/有重排序）的结果。}
  \label{fig:benchmark}
\end{figure*}

\subsection{基函数优化步数}

我们的方法通过递归细化构建 q 层层次结构。随着细化步骤数量的增加，粗层级 $\mathcal{M}^H$ 处基函数的支撑区域扩大，进而增加所得粗层级刚度矩阵 $\mathbf{K}^H$ 的密度。这种关系在 \Cref{fig:different-number-of-levels} 中针对三个不同模型进行了说明：螺母、体心立方（Body-Centered Cubic, BCC）微结构和三重周期性极小曲面（Triply Periodic Minimal Surface, TPMS）微结构。这些模型来自工业应用，在可视化中被裁剪以展示其内部结构。对于这些示例，我们使用与基准测试研究相同的边界条件求解静力平衡问题。

刚度矩阵中非零项的增加会降低共轭梯度（Conjugate Gradient, CG）方法的性能，因为其性能严重依赖于矩阵-向量乘积的效率。因此，我们的方法本质上涉及 CG 收敛所需迭代次数与每次迭代计算成本之间的权衡，如 \Cref{fig:different-number-of-levels} 所示。将这种权衡直接纳入基函数优化过程是复杂的，因为最优细化步数（q）因不同模型和问题配置而异。在其他结果中，默认选择 $q=40$ 层。

\begin{figure}[ht]
  \centering
  \includegraphics[width=\linewidth]{plot_figures/compare-steps/compare-steps.0524a.png}
  \caption{\textbf{求解线性变形问题时沿不同基函数优化步数的 CG 迭代比较。}不同模型间行为各异。刚度矩阵的非零密度随迭代次数减少而增加。时间成本在不同模型间获得不同的最小值。}
  \label{fig:different-number-of-levels}
\end{figure}

\subsection{压力测试}

\paragraph{病态单元引起的数值锁定}
我们将纸张镶嵌，中间部分像书一样粘合。当输入网格较差时，解可能落在基函数可表达空间之外。我们的方法在 $\mathcal{M}^{h+}$ 中扩展了这个空间，然后粗化到自由度（DoFs）更少的层级而不损失这种可表达性。我们展示了一个极端情况，其中网格因病态扁平单元而锁定（如 \Cref{fig:locking_issue} 所示）。自由度不足时，四面体无法弯曲，没有细化结果将保持在静止形状。使用我们的方法，纸张弯曲，而原始方法表现出伪影。延长到 $\mathcal{M}^{h+}$ 后，纸张自然弯曲。

\begin{figure}[ht]
  \centering
  \includegraphics[width=\linewidth]{figures/lock-zjx/lock.0523.png}
  \caption{一串病态单元相互连接。没有我们的方法，原始解被锁定且几乎为零（左）。使用我们方法中的细化步骤，我们摆脱了锁定问题但表现出大的拉伸而没有延长（中）。延长到细化网格 $\mathcal{M}^{h+}$ 后，我们得到合理的结果。}
  \label{fig:locking_issue}
\end{figure}

\paragraph{类四边形扁平单元}
类四边形扁平四面体，特别是 \Cref{fig:thin-elements} 中显示的第三个单元，对数值条件和解精度都构成重大挑战。这些单元经常出现在以六面体为主的重新网格化过程中，其中通过合并四面体生成六面体。虽然我们目前未实现混合单元网格的细化（将其留待未来工作），但我们的重点放在合并前的四面体网格上。\Cref{fig:quad-like-tet} 展示了一个类四边形单元高度分布的典型场景。我们的方法细化这些单元并使用优化基函数将其与邻居聚合。这种方法在共轭梯度方法中展示了优越的性能，无论是否使用 Jacobi 预条件子。

\begin{figure}[ht]
  \centering
  \includegraphics[width=\linewidth]{figpdf/nut.pdf}
  \caption{左上：本实验中的静力问题。右上：具有分布的类四边形单元和聚合的四面体网格。下：有无基函数优化的收敛比率比较。}
  \label{fig:quad-like-tet}
\end{figure}

\subsection{应用}

为了消除线性弹性的人工效应，
\begin{wrapfigure}[6]{r}{2cm}
    \vspace{-1em}
    \includegraphics[width=0.8\linewidth]{figpdf/warp.pdf}
\end{wrapfigure}
我们使用 \cite{HuangRS2011,pan_subspace_2015} 提出的旋转-应变坐标，配合泊松重建对结果进行后处理变形。该策略也在 \cite{chen2019material} 中采用。

\paragraph{切割仿真}
病态单元在切割仿真中频繁出现，如 \Cref{fig:cut-simulation} 所示。生成切割网格的一种常见且直接的方法是行进四面体方法，它根据顶点处的标量场值将每个四面体分割成多个新四面体。在此切割过程中，单元数量显著增加，导致病态四面体的激增（\Cref{p2:tab:statistics}），特别是在切割面附近。我们以准静态方式模拟此切割过程。使用基函数优化技术，刚度矩阵的最大特征值降低了一个数量级。这一优势将共轭梯度方法加速高达 $32\%$，而预计算成本仅为原始共轭梯度时间的 $1\%$。

\begin{figure}[t]
  \centering
  \includegraphics[width=0.8\linewidth]{figpdf/cut.pdf}
  \caption{\textbf{切割仿真。}病态四面体是切割过程的常见副产品。我们的方法有效识别对整体数值条件最有害影响的四面体，随后将其与相邻单元聚合。没有算子适应基函数优化，我们成功将系统的最大特征值降低一个数量级。这一显著改进带来了显著的好处：共轭梯度方法所需的计算时间减少了 $30\%$。}
  \label{fig:cut-simulation}
\end{figure}

\paragraph{动力学}
我们的基函数优化是一种预计算，使其特别适合线性系统被重复求解的动力学仿真。例如，我们将方法应用于包含大量类四边形病态四面体的螺母模型，使用向后欧拉积分，如 \Cref{fig:dynamics} 所示。我们的方法在使用大时间步时特别有益，其中刚度矩阵 $\mathbf{K}$ 主导线性系统 $M+h^2K$。在此示例中，我们将时间步设为 0.5。随着模拟帧数的增加，无论是否使用 Jacobi 预条件子，我们方法的优势都变得更加明显。

\begin{figure}[ht]
  \centering
  \includegraphics[width=\linewidth]{plot_figures/dynamics/dynamics-union.0524d.png}
  \caption{\textbf{隐式动力学仿真累积时间成本，有无基函数优化对比。}}
  \label{fig:dynamics}
\end{figure}
